\section{Methodology}
\label{sec:methodology}

Building on Martin and Mahoney's observation that heavy-tailed weight distributions indicate generalization quality, we design a measurement pipeline to test whether this theory extends to Vision Transformer attention mechanisms.

\subsection{Overview}

Our methodology has three components: (1) model collection from HuggingFace Hub, (2) heavy-tailed exponent computation via the Hill estimator, and (3) variance analysis across model populations.

\textbf{Design Rationale.} We use WeightWatcher, the canonical implementation of HT-SR analysis by Martin and Mahoney. This ensures methodological continuity with prior CNN analysis while testing on a new architecture class. By sampling from HuggingFace rather than a curated zoo, we test under realistic heterogeneity conditions.

\subsection{Model Collection}

We collect Vision Transformer models from HuggingFace Model Hub using the following criteria:

\begin{itemize}
    \item Filter: pipeline\_tag = image-classification
    \item Search: model name contains ``vit''
    \item Sort: downloads (descending)
    \item Target: 100 models
\end{itemize}

\textbf{Rationale.} Filtering by downloads prioritizes well-established models with community validation, reducing noise from abandoned or experimental checkpoints. The ViT architecture family provides a clean test case with standardized attention mechanisms.

\subsection{Heavy-Tailed Exponent Computation}

For each model, we compute the heavy-tailed exponent $\alpha$ using WeightWatcher. WeightWatcher performs singular value decomposition on each weight matrix and fits a power-law distribution $p(x) \propto x^{-\alpha}$ using the Hill estimator. The $\alpha$ value characterizes the tail behavior: lower values indicate heavier tails, correlating with better generalization per HT-SR theory.

We extract $\alpha$ values specifically for attention-related weight matrices (query, key, value projections) to isolate Transformer-specific behavior.

\subsection{Variance Analysis}

We compute two variance metrics:

\begin{enumerate}
    \item \textbf{Global variance} $\sigma_{\text{global}}$: Standard deviation of mean $\alpha$ across all models
    \item \textbf{Within-family variance} $\sigma_{\text{family}}$: Standard deviation within model families sharing training origin
\end{enumerate}

\textbf{Gate condition.} We set a threshold of $\sigma < 0.5$ for bounded variance, informed by prior CNN analysis where stable measurements enabled prediction.
